\chapter{Introduction}\label{ch:intro}

\paragraph{Motivation.} Elastomers filled with stiff particles combine
the compliance of the matrix with a tunable stiffness. For PDMS with glass
beads the contrast is extreme: the matrix has a shear modulus of order
$1\,$MPa and is nearly incompressible, while glass is about five orders of
magnitude stiffer. The macroscopic response is strongly nonlinear and depends
on the particle content. A concurrent FE$^2$ simulation resolves this by
solving a microstructural boundary-value problem at every macroscopic
integration point and load step~\cite{Kouznetsova2001,Miehe2002}. Its cost scales with the product of the macroscopic and microscopic problem
sizes. The alternative pursued here is
\emph{offline} homogenization. Representative volume elements (RVEs) are solved
once for many prescribed macroscopic deformations, and a surrogate learns the
homogenized energy so that it can act as a material law in the macroscopic
solver. For a Newton solver the surrogate must return the energy $\Psibar$
together with the first Piola stress $\Pbar$ and the consistent tangent $\Abar$,
and these three quantities must be mutually consistent.

\paragraph{Objective.} The project proposal (defended in SS~2026) defined the project as a
Python implementation of a thermodynamically consistent neural material law
inside NGSolve, built in two phases:
\begin{enumerate}
  \item \emph{Phase 1 (proof of concept):} a PyTorch--NGSolve material
        bridge, trained on a neo-Hookean law and verified on elementary FE problems;
  \item \emph{Phase 2 (multiscale):} periodic PDMS--glass RVEs,
        homogenized data generation, training of a volume-fraction-conditioned
        pICNN, and its deployment in a macroscopic FE simulation.
\end{enumerate}

\paragraph{Contributions and outcome.} All software components were
implemented and verified. An exactly rigid, freely rotating particle replaced the
proposed stiffness penalty, and a Schur-condensed energy Hessian replaced the
finite-difference RVE tangent. The verification covers the guideline's mandatory
element and single-element tests. A pressure/deviatoric split of the surrogate error
reveals and explains a large shear error, and an analysis shows that the
linear-element macroscopic solution is volumetrically locked.
\Cref{tab:proposal,tab:committee} compare the proposal milestones and the
committee's recommendations with their outcome.

\paragraph{Organization.} The chapters follow the order of the PPP
guideline: theory, numerical methods, available codes, flowchart, use of AI tools,
verification, assessment, computation time, conclusion, manual, literature and
git log.
